% ===========================================================================
%  THE MANUSCRIPT -- one body of text, shared by all three versions.
%  Never duplicate this file. Everything version-specific is a macro.
% ===========================================================================

\begin{abstract}
Collaborative revision of a LaTeX manuscript usually degenerates into a folder
of files named \texttt{v3\_final\_JD\_comments\_REALfinal.tex}. This template
keeps a single source and derives three documents from it: the working master,
the comment-driven version circulated among collaborators, and the clean file
that goes to the journal.
\end{abstract}

\section{Introduction}

The problem is not that people disagree about a manuscript; it is that the
disagreement has nowhere to live. A comment written in the margin of a PDF is
detached from the sentence it concerns, and a sentence rewritten in a private
copy is detached from the reason it was rewritten.
\comment{co1}{Is this opening too polemical for a methods paper? I would soften
it.}%
\reply{sam}{I would rather keep it. The paper is short and the framing pays for
itself in the discussion.}

Here the comment and the text are the same object. Every collaborator writes
into the same file, and the file knows how to hide what a given reader should
not see.
\note{sam}{Private reminder: check whether the journal allows line numbers in
the submitted version.}

\section{How the markup behaves}

\subsection{Comments and responses}

A comment is raised with \verb|\comment{who}{text}| and answered with
\verb|\reply{who}{text}| when nothing changes, or with
\verb|\resolved{who}{text}| when something did.
\comment{rev}{The description of the estimator in Section~2 is too terse to
reproduce.}%
\resolved{sam}{Expanded into a numbered procedure with the tolerance stated
explicitly.}

Inline comments keep the working discussion beside the affected text. The
formal editor-facing wording is maintained separately in
\texttt{point-by-point.tex}, where the authors can polish each response.

\subsection{Changes and suggestions}

Text written specifically to address a comment is wrapped in
\verb|\changed{who}{...}|. It is
\changed{sam}{highlighted in the author's colour while the manuscript is under
review, and becomes ordinary text in the submitted version.}

A suggestion is different: it is a proposal that has not been acted on yet.
\verb|\suggest{who}{old}{new}| shows both readings side by side while the
document is being reviewed, and collapses to one of them in the submitted
version according to the accept policy. The convergence rate is
\suggest{co2}{roughly linear}{linear in the number of iterations, with a
constant that depends on the conditioning of $A$}.

Insertions and deletions work the same way.
\suggestadd{co2}{We assume throughout that $A$ is symmetric positive definite.}
The next sentence is redundant and has been proposed for removal.
\suggestdel{co1}{As is well known, this assumption is standard in the
literature.}

A single suggestion can override the global policy, which is how a co-author
records a proposal that was considered and turned down:
\suggest[reject]{co1}{the residual}{the error}. In the submitted version this
prints \emph{the residual}, whatever the accept policy says, and the record of
the discussion survives in the comment-driven version.

\subsection{Mathematics is not a special case}

The markup is inline text, so it lives happily next to displayed equations:
\begin{equation}
  \label{eq:fixedpoint}
  x_{k+1} = x_k - \alpha_k \, A^{-1} r_k ,
  \qquad r_k = b - A x_k .
\end{equation}
Equation~\eqref{eq:fixedpoint} converges for
\suggest{co2}{$0 < \alpha_k < 1$}{$0 < \alpha_k < 2/\lambda_{\max}(A)$}.%
\comment{co2}{The step-size condition as written is sufficient but far from
sharp.}
\revtodo{sam}{Add the sharp constant and a one-line proof in the appendix.}

\section{Results}

\begin{table}[htbp]
  \centering
  \caption{Markup inside floats behaves like markup anywhere else.}
  \label{tab:demo}
  \begin{tabular}{lrr}
    \toprule
    Method & Iterations & Time (s) \\
    \midrule
    Baseline   & 412 & 3.81 \\
    Proposed   & \suggest{co1}{97}{94} & 0.92 \\
    \bottomrule
  \end{tabular}
\end{table}

Table~\ref{tab:demo} reports the iteration counts.
\comment{co1}{These numbers are from the old run. Please regenerate against the
current code before we circulate this.}%
\resolved{sam}{Regenerated; the proposed method now needs 94 iterations.}

\masteronly{%
\subsection*{Working scratch (master only)}
This subsection exists only in \texttt{main-original.tex}. It is where the
derivations that nobody else needs to read can sit without cluttering the
circulated version.}

\finalonly{%
\subsection*{Funding}
This work was supported by \dots{} This paragraph appears only in the submitted
version, so it does not distract reviewers of the working drafts.}

\section{Conclusion}

One manuscript source, three document views, and one explicit response-letter
source. The alternative is a folder of files nobody trusts.
\comment{co2}{Add a sentence on how this compares with the \texttt{changes}
package.}
